\section{Orientation and Scope}
\label{sec:orientation}

\begin{roadmapbox}
This prototype has a narrower objective than a final textbook. It is designed
to prove that the theory can support a thicker table of contents, repeatable
case analysis, and learner-facing exercises. In business terms, this chapter
defines the product specification for the manuscript: what the book explains,
what it refuses to explain, and why the model is useful.
\end{roadmapbox}

\subsection{The Core Question}

Why do agents often fail to initiate a beneficial activity even when the
activity is clearly endorsed by their reflective judgment? The model proposed
here answers that question by separating three objects:
\begin{enumerate}[nosep]
  \item the current behavioural state,
  \item the barrier that must be crossed to leave that state,
  \item the timing conditions under which volitional effort has the highest
        practical leverage.
\end{enumerate}

\subsection{What This Book Is Trying To Deliver}

The manuscript aims to do four jobs:
\begin{enumerate}
  \item provide a compact vocabulary for talking about initiative failure,
  \item show why planning is structurally necessary rather than morally
        optional,
  \item organize real-life cases into a reusable analytical framework,
  \item supply exercises so that the reader can use the framework rather than
        merely admire it.
\end{enumerate}

\subsection{What This Book Does Not Claim}

\begin{assumption}[Limited ambition]
Framework Inertia Theory is a semi-formal decision model. It does not claim to
replace neuroscience, clinical psychology, or a full theory of action.
\end{assumption}

\begin{remark}
This boundary matters. A textbook becomes more credible when it names its
coverage limits early. The model here is intended to be operational: it should
help a learner diagnose blocked transitions and design plans that survive
ordinary depletion.
\end{remark}

\subsection{Chapter Map}

\begin{longtable}{@{}p{0.16\textwidth} p{0.74\textwidth}@{}}
\toprule
\textbf{Chapter} & \textbf{Function} \\
\midrule
1 & Define the formal objects: states, inertia, barriers, and transition load. \\
2 & Explain why free will has variable leverage and identify decision windows. \\
3 & Show why planning is a structural intervention, not a decorative habit. \\
4 & Stress-test the framework across study, exercise, digital distraction, and recovery cases. \\
5 & Translate theory into implementation patterns for individuals and systems. \\
6 & Provide exercises, prompts, and a self-audit worksheet for applied use. \\
Appendix A & Stabilize notation and provide templates that can be reused in later drafting rounds. \\
\bottomrule
\end{longtable}

\begin{summarybox}
The commercial value of a richer textbook starts here: a learner can see where
the book is going, how each chapter contributes, and why the model is meant to
be used as an operating system rather than a motivational speech.
\end{summarybox}

